\chapter{Operating the Stack: A Debugging Field Manual}
\label{ch:operations}

This chapter is the one to reread with a broken cluster in front of
you. It assembles the debugging workflow the earlier chapters earned,
ordered the way real incidents unfold.

\section{The four questions}

Almost every Kubernetes incident is located by asking, in order:

\begin{enumerate}
  \item \textbf{Is the pod running?} \code{kubectl -n ts get pods} ---
  read the STATUS and RESTARTS columns; the taxonomy below.
  \item \textbf{What does the platform say about it?}
  \code{kubectl -n ts describe pod <pod>} --- the Events section at the
  bottom is the kubelet and scheduler narrating their decisions:
  image pulls, probe failures, OOM kills, scheduling refusals.
  \item \textbf{What does the app say?} \code{kubectl -n ts logs
  <pod>} (add \code{--previous} for the container that just died ---
  the single most-forgotten flag in incident response).
  \item \textbf{Can I reproduce its view of the world?}
  \code{kubectl -n ts exec -it <pod> -- sh} and try the connection the
  app claims is failing.
\end{enumerate}

\section{A status taxonomy}

\begin{table}[h]
\centering\small
\begin{tabular}{@{}lp{6.6cm}@{}}
\toprule
Status & Meaning and first move \\
\midrule
\code{Pending} & unschedulable: requests exceed free
resources, or PVC unbound. \code{describe} shows which. \\
\code{ImagePullBackOff} & registry unreachable, name typo'd, or the
HTTPS-vs-HTTP trap of Chapter~13. \\
\code{CrashLoopBackOff} & process exits soon after start; backoff
doubles between retries. \code{logs --previous}. \\
\code{OOMKilled} (exit 137) & memory limit hit; no Java trace exists
--- Chapter~17's kernel kill. Raise limit or shrink heap. \\
\code{Running 0/1} & alive but readiness failing: normal during
startup; a problem when persistent --- probe the probe's own URL. \\
\code{Completed}/\code{Error} & for Jobs; for a service pod,
\code{Error} means the process exited nonzero once. \\
\bottomrule
\end{tabular}
\caption{What the STATUS column is trying to tell you.}
\end{table}

\begin{handson}[break it on purpose --- three drills]
Deliberate sabotage, on a quiet day, teaches more than any outage:
\begin{lstlisting}[language=bash]
# 1. wrong image tag -> watch ImagePullBackOff, then roll back
kubectl -n ts set image deployment/course-service \
    course-service=localhost:5000/course-service:nope
kubectl -n ts rollout undo deployment/course-service

# 2. scale the database down -> watch dependents degrade honestly
kubectl -n ts scale statefulset/postgres-course --replicas=0
kubectl -n ts get pods -w     # course-service readiness reacts
kubectl -n ts scale statefulset/postgres-course --replicas=1

# 3. kill Eureka -> verify existing routes survive on client caches
kubectl -n ts delete pod -l app=discovery-server
\end{lstlisting}
Each drill pairs a chapter's theory with its observable behavior:
13 (naming), 17--18 (probes and endpoints), 3 (client-side caching).
\end{handson}

\section{Reading a service's world from inside}

When logs claim a connection failure, verify the claim from the pod's
own network position --- DNS and reachability differ between your
machine and the pod:

\begin{lstlisting}[language=bash]
kubectl -n ts exec -it deploy/course-service -- sh
wget -qO- http://config-server:8888/actuator/health   # DNS + service up?
nc -zv postgres-course 5432                           # TCP reachable?
env | grep SPRING                                     # what was injected?
\end{lstlisting}

That last command settles arguments fast: it prints what configuration
the pod \emph{actually} received, which beats re-reading manifests.

\section{The routine that keeps weekends free}

\begin{description}
  \item[After every deploy] the pipeline already verifies rollouts;
  glance at \code{kubectl -n ts get pods} anyway --- RESTARTS creeping
  upward is tomorrow's outage announcing itself.
  \item[Weekly] \code{kubectl top pods -n ts} (requires
  metrics-server, bundled with k3s): compare actual usage against the
  requests of Chapter~17 and correct the fiction. Check disk:
  \code{df -h} on the server --- registries and log files grow;
  \code{registry:2} needs occasional garbage collection.
  \item[Before demos] \emph{do not build} --- Chapter~5's hardware box:
  two cores mean a Maven build and a demo contend for the same CPU
  that keeps readiness probes answering.
\end{description}

\begin{hardware}
ufw on the server must exempt the pod and service CIDRs
(\code{10.42.0.0/16}, \code{10.43.0.0/16}) --- firewalling them breaks
DNS and pod-to-pod traffic in ways that impersonate application bugs:
intermittent timeouts, resolvable-sometimes names. When \emph{weird}
networking happens on k3s, suspect the host firewall before the app.
And k3s reads \code{registries.yaml} only at startup --- after editing,
\code{systemctl restart k3s}, which briefly restarts every pod: not
during a demo.
\end{hardware}

\begin{pitfall}[debugging the wrong layer]
The most expensive habit in Kubernetes debugging is starting where the
symptom appeared instead of where the cause lives. A browser 502 can be
Traefik (no endpoints), the gateway (Eureka staleness), the service
(readiness failing), the database (pod evicted), or the disk (full).
The four questions exist to walk the layers in order; resist the
temptation to open application logs first just because they are the
most familiar layer. \code{describe} before \code{logs} --- the
platform usually already knows.
\end{pitfall}

\begin{quiz}
\begin{enumerate}
  \item Order the four questions and name the command for each. Which
  one uses a flag people forget, and when does that flag matter?
  \item \code{CrashLoopBackOff} vs.\ \code{Running 0/1}: different
  layers are failing --- which probe concept distinguishes them?
  \item A service claims it cannot reach Postgres; from your laptop
  Postgres answers. What is the next diagnostic step and why is your
  laptop's success irrelevant?
  \item Why can a host firewall change masquerade as an application
  bug on k3s, and which two CIDRs must stay exempt?
\end{enumerate}
\end{quiz}
